\section{Query Processing and Optimization}

What is query optimization? Why is it necessary?
query optmization is necessary when stuff is very slow, which is especially needed for OLAP applications bc blah
blah


What query optimization strategies are available?
    Correct database design and normalization
    star schema allows for signifcicantly faster queries, as a lot of the data is already aggregated into a few tables
    which results in less joins and faster access to data

    Use of indexes – primary and secondary indexes
    indexing is good for use on frequently accessed columns that aren't already primary or foreign keys
    keep in mind write trade off for OLTP, and that indexes dont always help

    Query restructuring
    Another effective method used to optimize queries is the use of commmon table expressions (CTEs), specifically
    the MATERIALIZED view that stores the results of a subquery in physical memory, which prevents potential re-computation that
    signifcicantly increases a query's runtiem. This is best demonstrated in the company monthly ranking query which goes from 
    approx 280k ms execution time down to around 3-4000ms. Query 2 also utilizes this but to a lesser degree, cutting runtime in about
    half as it only needds to compute for *cur* once rather than twice

    Add baseline vs optimized query screenshots for 1 and 2

Provide illustrations – diagrams, SQL codes, schema – to illustrate the optimization strategy.


Cite related literature to support your design decisions  